% SOURCE_AUTHORITY: sga5_fr_workpass.tex, lines 9248--9290 (LF-normalized unit).
% SOURCE_SHA256: 791F4EFFC5E02832D5D77ED03518C8156D6F07E4C8238B03545DB93D883FBB28
\subsubsection*{2.2}
Sean $S$ un esquema y $E$ un $\mathcal O_S$-módulo localmente libre de rango $r+1$. Se denota por
\[
P=P_S(E)\xrightarrow{p}S
\]
el fibrado proyectivo correspondiente. El $\mathcal O_P$-módulo canónico $\mathcal O_P(1)$ define, como es sabido, gracias a la sucesión exacta de Kummer, un elemento
\[
\xi\in H^2(P,\mu)
\]
y se denotará por $\eta$ la imagen de $\xi$ por el morfismo canónico
\[
H^2(P,\mu)\longrightarrow H^0(S,R^2p_*(\mu))
\]
procedente de la sucesión espectral de Leray
\[
H^i(S,R^jp_*(\mu))\Rightarrow H^{i+j}(P,\mu).
\]
La clase $\xi$ se identifica con una flecha $A_P\to\mu_P[2]$ de $D(P)$, que define por transposición $\mu_P^{\otimes -1}[-2]\to A_P$, de donde, por la adjunción usual, se obtiene una flecha
\[
c:\mu_S^{\otimes -1}[-2]\longrightarrow Rp_*(A_P),
\]
que coincide con $\eta$ en la cohomología de grado $2$.

Por producto tensorial, la flecha $c$ define, gracias a (2.2.1), unas flechas
\[
c_i:\mu_S^{\otimes -i}[-2i]\longrightarrow Rp_*(A_P),
\]
para todo entero $i\ge1$, y se conviene en denotar por $c_0$ la flecha de adjunción
\[
A_S\longrightarrow Rp_*(A_P).
\]
Se define así una flecha
\[
\gamma=\bigoplus_{i=0}^{r}c_i:
\bigoplus_{i=0}^{r}\mu_S^{\otimes -i}[-2i]\longrightarrow Rp_*(A_P)
\]
de $D(S)$.

\begin{statement}{Teorema 2.2.1}
La flecha $\gamma$ es un isomorfismo.
\end{statement}

Una vez definida la flecha $\gamma$, (2.2.1) equivale al corolario siguiente.
